\documentclass{../template/note}

\author{Oliver West}
\title{Groups, rings and modules}
\date{}

\begin{document}

    \maketitle

    \tableofcontents

    \lecture{Lecture 1 -- 22/01/26}

    \section{Groups -- a reminder}

    \begin{definition}
        A \emph{group} is a triple $(G, \cdot, e)$ where $G$ is a set, 
        \begin{align*}
            \cdot: G \times G \to G
        \end{align*}
        is a function, and $e \in G$ such that
        \begin{itemize}
            \item \emph{Associativity}. $(a \cdot b) \cdot c = a \cdot (b \cdot c)$;
            \item \emph{Identity}. For all $a \in G$, $a \cdot e = e \cdot a = a$;
            \item \emph{Inverses}. For all $a \in G$, there exists $a^{-1} \in G$ such that $a \cdot a^{-1} = a^{-1} \cdot a = e$.
        \end{itemize}
        The size of $G$, $|G|$, is called the \emph{order} of the group. 
    \end{definition}
    
    \begin{proposition}[Uniqueness of inverses]
        Inverses are unique.
    \end{proposition}

    \begin{proof}
        Left as an exercise (see 1A).
    \end{proof}
    
    \begin{definition}[Subgroups]
        If $G$ is a group and $H \subseteq G$, then we say $H$ is a \emph{subgroup} if
        \begin{itemize}
            \item $e \in H$;
            \item for all $a$, $b \in H$, $a \cdot b \in H$;
            \item $(H, \cdot, e)$ is a group.
        \end{itemize}
        In order to say item 3, we must first establish 1 and 2.
    \end{definition}
    
    \begin{proposition}
        A non-empty subset $H \subseteq G$ is a subgroup if and only if, for all $h_1$, $h_2 \in H$, $h_1h_2^{-1} \in H$.
    \end{proposition}
    
    \begin{proof}
        Left as an exercise (see 1A).
    \end{proof}
    
    \begin{definition}
        A group $G$ is called \emph{abelian} if, for all $g_1$, $g_2 \in G$, $g_1g_2 = g_2g_1$.
    \end{definition}
    
    \begin{example}
        \begin{itemize}
            \item We have $(\ZZ, +)$, $(\QQ, +)$, $\RR$, $\CC$, $\ZZ/n\ZZ$.
            \item The symmetric groups $S_n$: permutations of $\{1, \dots, n\}$ under composition. More generally for any set $X$, we have $\Sym(X)$.
            \item The dihedral groups $D_{2n}$: the group of symmetries of a regular $n$-gon.
            \item The general linear group over $\RR$, $\GL(n, \RR)$ or $\GL_n(\RR)$, is the set of $\{f: \RR^n \to \RR^n \mid f \text{ is an }\RR\text{-linear bijection}\}$. Since $\GL(n, \RR)$ is made of functions, we can think of it a subset of $\Sym(\RR^n)$. (Equivalently, it's the set of $n$ by $n$ invertible matrices with elements in $\RR$).
        \end{itemize}
        Some examples of subgroups:
        \begin{itemize}
            \item The even permutations $A_n \leq S_n$ (the alternating group);
            \item Rotational symmetries of an $n$-goin, $C_n \leq D_{2n}$;
            \item Determinant one matrices $\SL(n, \RR) \leq \GL(n, \RR)$.
        \end{itemize}
        Other examples include $C_2 \times C_2$, $Q_8 = \{\pm 1, \pm i, \pm j, \pm k\}$.        
    \end{example}
    
    \begin{definition}
        Let $G$ be a group, $g \in G$ and $H \leq G$ a subgroup. Then the \emph{left coset}
        \begin{align*}
            gH = \{gh: h \in H\}
        \end{align*}
    \end{definition}
    Basic properties include that $g \in gH$, so cosets are non-empty. Also, cosets partition $G$ -- they have no non-trivial overlaps and every $g$ lies in some coset. Furthermore, the map
    \begin{align*}
        H &\to gH \\
        h &\mapsto gh
    \end{align*}
    defines a bijection.
    
    \begin{theorem}[Lagrange]
        If $G$ is a finite group and $H \leq G$ then
        \begin{align*}
            |G| = |H| \cdot |G:H|
        \end{align*}
        where $|G:H|$, the index of $H$ in $G$, is the number of cosets of $H$ in $|G|$.
    \end{theorem}
    
    \begin{definition}
        Let $G$ be a group and let $g \in G$. The \emph{order} of $g$, $\ord(g)$, is the smallest positive integer $n$ such that $g^n = e$, or infinity if no such $n$ exists.
    \end{definition}
    
    \begin{proposition}
        Let $g \in G$. Then, $\ord(g) \mid |G|$.
    \end{proposition}
    
    \begin{proof}
        Let $g \in G$ and consider $H = \{e, g, g^2, \dots, g^{n-1}\}$ where $n = \ord(g)$. This is a subgroup: it is non-empty, and $g^r \cdot g^{-s} = g^{r-s} \in H$ (after potentially multiplying by $g^n = e$). Also, the elements of $H$ are distinct -- if $g^i = g^j$ for $0 \leq i < j \leq n$, then $g^{j-i} = e$, contradicting the minimality of $\ord(g)$. Hence $|H| = ord(g)$ -- now apply Lagrange.
    \end{proof}
    
    \section{Normal subgroups, quotients and homomorphisms}

    Let $H \leq G$. When are two cosets of $H$ in $G$ equal? If $gH = g'H$, then $g \in g'H$, so $g = g'h$ for some $h \in H$. In other words, $gH = g'H \iff (g')^{-1}g \in H$.

    We denote the set of left cosets of $H \leq G$ by $G/H = \{gH: g \in G\}$. Does $G/H$ carry a natural group structure?

    \lecture{Lecture 2 -- 24/01/26}

    We propose the `natural' group structure:
    \begin{align*}
        g_1H \cdot g_2H = (g_1g_2)H
    \end{align*}
    
    But this rule may depend on the coset representatives we pick. We must check that it doesn't. Suppose $g_2H = g_2'H$, so $g_2' = g_2h$ for some $h \in H$. Then
    \begin{align*}
        g_1H \cdot g_2'H = (g_1g_2')H = g_1g_2hH = g_1g_2H
    \end{align*}
    
    On the other hand, if $g_1H = g_1'H$, so $g_1' = g_1h$ for $h \in H$. Then
    \begin{align*}
        g_1'H \cdot g_2H = g_1hg_2H
    \end{align*}
    
    and that is not automatically equal to $g_1g_2H$ -- we need $g_2^{-1}hg_2 \in H$ to hold for all $g_2 \in G$, $h \in H$.

    \begin{definition}[Normal subgroups]
        A subgroup $H \leq G$ is \emph{normal} if, for all $g \in G$, $h \in H$, $ghg^{-1} \in H$. We write this as $H \norm G$.
    \end{definition}
    
    \begin{remark}
        When $H \norm G$, then $G/H$ inherits a group structure by the rule above. We will refer to this as the \emph{quotient group of $G$ by $H$}.
    \end{remark}
    
    \begin{definition}[Homomorphisms]
        Let $G$ and $H$ be groups. A homomorphism from $G$ to $H$ is a function $f: G \to H$ such that, for all $g_1$, $g_2 \in G$, $f(g_1 g_2) = f(g_1) f(g_2)$.
    \end{definition}
    
    Results such as $f(e_G) = e_H$ follow immediately.
    
    \begin{lemma}
        If $f: G \to H$ is a homomorphism, then $f(g^{-1}) = f(g)^{-1}$.
    \end{lemma}
    
    \begin{proof}
        Calculate $f(g \cdot g^{-1})$ in two ways.
        \begin{enumerate}[label=\protect\circled{\arabic*}]
            \item $f(gg^{-1}) = f(e_G) = e_H$;
            \item $f(gg^{-1}) = f(g) f(g^{-1})$
        \end{enumerate}
        so $f(g) f(g^{-1}) = e_H$ and the result follows.        
    \end{proof}
    
    \begin{definition}
        Let $f: G \to H$ be a homomorphism. The \emph{kernel} of $f$ is
        \begin{align*}
            \ker f = \{g \in G : f(g) = e_H\}
        \end{align*}
        and the \emph{image} of $f$ is
        \begin{align*}
            \Im f = \{h \in H: h = f(g) \text{ for some } g\}
        \end{align*}     
    \end{definition}
    
    \begin{proposition}
        Let $f: G \to H$ be a homomorphism. Then $\ker f \norm G$ and $\Im f \leq H$.
    \end{proposition}
    
    \begin{proof}
        Observe $e \in \ker f$ so $\ker f$ non-empty. Now, if $x, y \in \ker f$, then
        \begin{align*}
            f(xy^{-1}) = f(x)f(y)^{-1} = e
        \end{align*}
        so $xy^{-1} \in \ker f$. For normality, let $x \in G$ and $g \in \ker f$. Now
        \begin{align*}
            f(xgx^{-1}) = f(x) e f(x)^{-1} = e
        \end{align*}
        so $xgx^{-1} \in \ker f$ and hence $\ker f \norm G$.
        
        For the image, observe that, if $a$, $b \in \Im f$ with $a = f(x)$, $b = f(y)$, then $ab^{-1} = f(xy^{-1}) \in \Im f$.
    \end{proof}
    
    \begin{remark*}
        Observe $f: G \to G/H$ taking $g \mapsto gH$ is a homomorphism. This is a nice way of expressing its `naturalness'.
    \end{remark*}
    
    \begin{definition}
        A homomorphism $f: G \to H$ is an \emph{isomorphism} if it is a bijection. Two groups $G$ and $H$ are said to be \emph{isomorphic}, $G \cong H$, if there exists an isomorphism between them.
    \end{definition}
    
    \begin{theorem}[First isomorphism theorem]
        Let $f: G \to H$ be a homomorphism, then $\ker f$ is normal in $G$ and the homomorphism
        \begin{align*}
            G/\ker f &\to \Im f \\
            g \ker f &\mapsto f(g)
        \end{align*}
        is an isomorphism. In particular, $G/\ker f \cong \Im f$.                
    \end{theorem}
    
    \begin{proof}
        We have already shown the normality result. Consider the map
        \begin{align*}
            \phi: G/\ker f &\to \Im f \\
            g\ker f &\mapsto f(g)
        \end{align*}
        \begin{itemize}
            \item \emph{$\phi$ is well-defined}. If $g \ker f = g'\ker f$, then $g'g^{-1} \in \ker f$, so $f(g'g^{-1})= e$ which gives $f(g) = f(g')$ as required.
            \item \emph{$\phi$ is surjective}. Let $h \in \Im f$. Then $h = f(g)$ for some $g \in G$, and so $h = \phi(g \ker f)$.
            \item \emph{$\phi$ is injective}. Suppose $\phi(g \ker f) = \phi(g' \ker f)$, then $f(g) = f(g')$. But then $g'g^{-1} \in \ker f$, so $g \ker f = g' \ker f$.
        \end{itemize}
    \end{proof}
    
    \begin{example*}
        Consider the homomorphism
        \begin{align*}
            \exp: (\CC, +) &\to (\CC^*, \times) \\
            z &\mapsto e^z
        \end{align*}
        Then $\ker \exp = 2\pi i\ZZ$, so $\CC^* \cong \CC/2\pi i \ZZ$.
    \end{example*}
    
    \begin{theorem}[Second isomorphism theorem]
        Let $H \leq K$ and $K \norm G$. Then
        \begin{align*}
            HK = \{hk : h \in H, k \in K\}
        \end{align*}
        is a subgroup of $G$. Furthermore, $H \cap K$ is normal in $H$, and there is an isomorphism $HK/K \cong H/H \cap K$.
    \end{theorem}
    
    \begin{proof}
        We will show the results in order.
        \begin{itemize}
            \item \emph{$HK$ is a subgroup}. It contains $e_G$, so is non-empty. Now consider $x = hk$ and $y = h'k'$. We'll show $yx^{-1} \in HK$. Observe
            \begin{align*}
                yx^{-1} &= h'k'(hk)^{-1} \\
                &= h'k'k^{-1}h^{-1} \\
                &= (h'h^{-1})h(k'k^{-1})h^{-1} \\
                &= h'h^{-1}k''
            \end{align*}
            by normality of $K$.
            \item \emph{$H \cap K$ is a kernel}. Consider
            \begin{align*}
                \phi: H &\to G/K \\
                h &\mapsto hK
            \end{align*}
            by restricitng the domain from $G \to G/K$ -- the kernel is $H \cap K$.
            \item \emph{Final statement}. By the first isomorphism theorem, we have
            \begin{align*}
                \Im \phi \cong H/(H \cap K)
            \end{align*}
            so we must show $\Im \phi \cong HK/K$. But $\Im \phi$ is the set of $K$-cosets in $G$ that are representable by elements of $H$, i.e. $\{hK : h \in H\} \subseteq G/K = HK/K$.
        \end{itemize}
    \end{proof}
    
    \subsection{Correspondence between subgroups of $G$ and $G/K$}

    Let $p: G \to G/K$ be the quotient homomorphism. We have a bijection
    \begin{align*}
        \{\text{subgroups of } G \text{ containing } K\}& \leftrightarrow \{\text{subgroups of } G/K\} \\
        K \norm L \leq G &\mapsto L/K \leq G/K \\
        \{g \in G : gK \in A\} &\mapsfrom A \leq G/K
    \end{align*}
    
    The forward bijection is just applying the projection map on a subset, and, since the quotient map is a homomorphism, its image is a subgroup; and the backward bijection is actually just applying the preimage of $p$. It is an exercise to check that the forward and backward directions are inverses.

    \begin{remark*}
        In fact the correspondence matches normal subgroups on the two sides.
    \end{remark*}
    
    \begin{theorem}[Third isomorphism theorem]
        Let $K \leq L \leq G$ be normal subgroups of $G$. Then there exists an isomorphism $(G/K)/(L/K) \cong G/L$.
    \end{theorem}
    
    \begin{proof}
        Define a map
        \begin{align*}
            \phi: G/K &\to G/L \\
            gK &\mapsto gL
        \end{align*}
        \begin{itemize}
            \item \emph{$\phi$ is well-defined}. If $gK = g'K$, then $g'g^{-1} \in K \subseteq L \implies gL = g'L$.
            \item \emph{$\phi$ is a homomorphism}. Clear.
        \end{itemize}
        We have $\ker \phi = \{gK: gL = L\} = \{gK: g \in L\} = L/K$. Hence we are done by the first isomorphism theorem.
    \end{proof}
    
    \begin{definition}[Simplicity]
        A group $G$ is called \emph{simple} if its only normal subgroups are $G$ and $\{e\}$.
    \end{definition}
    
    \begin{proposition}
        Let $G$ be an abelian group. Then $G$ is simple if and only if $G \cong C_p$ for some prime $p$.
    \end{proposition}
    
    \begin{proof}
        If $G \cong C_p$, then any non-identity element generates $G$ by, for example, Lagrange. So $G$ is simple.

        If $G$ is simple and abelian, then take any $g \neq e \in G$ and consider $\gen g$. This is a subgroup, and since $G$ is abelian, it is automatically normal. Hence, by simplicity, $\gen g = G$, and so $G$ is cyclic.

        If $G$ is infinite and cyclic, then $G \cong \ZZ$, but $\ZZ$ is not simple, so this is a contradiction. Hence $G$ is finite and $G \cong C_m$. If $q$ divides $m$, take the subgroup generated by $g^{m/q}$, where $g$ generates $G$. Normality is automatic, so $q=1$ or $q=m$ by simplicity -- in other words, $m$ is prime.
    \end{proof}
    
    \begin{theorem}
        Let $G$ be a finite group. Then, there exist subgroups
        \begin{align*}
            G = H_1 \triangleright H_2 \triangleright H_3 \triangleright \cdots \triangleright H_n = \{e\}
        \end{align*}
        such that $H_i/H_{i+1}$ is simple for all $i$.        
    \end{theorem}
    
    \begin{proof}
        If $G$ is simple, $H_2 = \{e\}$, and we're done. Otherwise, let $H_2$ be a proper normal subgroup of maximal order in $G$. We claim $G/H_2$ is simple. If not, consider the quotient map $\pi: G \to G/H_2$.

        If $N \norm G/H_2$ is a proper normal subgroup, by the correspondence theorem, we can find a proper normal subgroup in $G$ containing $H_2$. But this contradicts the maximality of $H_2$.

        Now we can just iterate -- note that $|H_{i+1}| < |H_i|$, so we will eventually terminate.
    \end{proof}
    
    \begin{remark*}
        This sequence is not necessarily unique.
    \end{remark*}
    
    \section{Group actions and permutations}

    \begin{definition}
        Let $X$ be a set. Let $\Sym(X)$ be the set of permutations on $X$. The identity is the identity function, and the group law is given by composition. If $X = \{1, \dots, n\}$, then we write $\Sym(X) = S_n$.
    \end{definition}
    
    \begin{itemize}
        \item It's convenient to write $\sigma \in S_n$ as a product of disjoint cycles.
        \item Every $\sigma \in S_n$ is a product of transpositions, and the number of such transpositions gives rise to the \emph{sign} of $\sigma$.
        \begin{align*}
            \sign: S_n &\to \{\pm 1\} \\
            \sigma &\mapsto \begin{cases}
                +1 & \sigma \text{ is a product of an even number of transpositions} \\
                -1 & \sigma \text{ ... odd ...}
            \end{cases}
        \end{align*}
        $\sign$ is a homomorphism, and is surjective when $n \geq 3$.
        
    \end{itemize}
    
    
    
    
    
    
    
    
\end{document}
